\chapter{力的功}\label{ch:g11:work-of-force}

把衣柜推过房间，你就做了功；顶着墙推
一整个下午，不论肌肉感觉如何，物理上都计为
零。功是力沿\emph{运动}所输送的能量——
即传递中的能量（见\cref{ch:g10:energy-conservation}），以焦耳计量。本章
计算\cref{ch:g11:forces-and-motion}中的力——
斜向拉力、重力、摩擦力，并给出做功的快慢：功率。

\section{恒力的功}

\begin{definition}[恒力的功]\label{def:g11:work-of-force:work}
恒力 $\vect F$ 作用于沿直线段
从 $A$ 运动到 $B$ 的物体，所做的\emph{功}\index{功（力的）}为
\[
W_{AB}(\vect F) = F \times AB \times \cos\theta ,
\]
其中 $F$ 为力的大小（单位为 \unit{N}），$AB$ 为
位移的长度（单位为 \unit{m}），$\theta$ 为
$\vect F$ 与 $\vect{AB}$ 的夹角。功的单位是
\emph{焦耳}\index{焦耳}（\unit{J}）：一焦耳是一牛顿的力
沿其自身方向移动一米所做的功。
\end{definition}

\begin{remark}\label{rem:g11:work-of-force:scal}
$F \times AB \times \cos\theta$ 恰是本学年数学教材中的
数量积 $\vect F \cdot \vect{AB}$。只有
\emph{沿}运动的分量 $F\cos\theta$ 才做功。
\end{remark}

\begin{omfigure}
\begin{tikzpicture}[font=\small]
  \draw (-0.4,0) -- (6.4,0);
  \draw[fill=omExo!15] (0.4,0) rectangle (1.5,0.7);
  \coordinate (P) at (0.95,0.35);
  \draw[dashed, omExo] ($(P)+(2.29,1.61)$) -- ($(P)+(2.29,0)$);
  \draw[-Stealth, omDef, very thick] (P) -- ($(P)+(35:2.8)$)
    node[above] {$\vect F$};
  \draw[-Stealth, omMeth, thick] (P) -- ++(2.29,0)
    node[pos=0.75, below] {$F\cos\theta$};
  \draw (P) ++(1.1,0) arc[start angle=0, end angle=35, radius=1.1];
  \node at ($(P)+(17:1.38)$) {$\theta$};
  \draw[-Stealth, omThm, thick] (0.95,-0.7) -- (4.95,-0.7)
    node[midway, below] {$\vect{AB}$};
  \fill (0.95,-0.7) circle (1.2pt) node[below left] {$A$};
  \fill (4.95,-0.7) circle (1.2pt) node[below right] {$B$};
\end{tikzpicture}

{\small 只有力在位移方向上的投影 $F\cos\theta$
做功：$W_{AB}(\vect F) = F \times AB \times \cos\theta$。}
\end{omfigure}

\begin{example}[拖着的行李箱]\label{ex:g11:work-of-force:suitcase}
旅客在航站楼内拖着行李箱移动 \qty{300}{m}，
拉带受拉力 $F = \qty{40}{N}$，与地面的夹角为 $\theta = \qty{50}{\degree}$，
则 $W = 40 \times 300 \times \cos\qty{50}{\degree} \approx
\qty{7.7}{kJ}$。若拉带水平（$\theta = 0$），
同样的力将做功 \qty{12}{kJ}：约三分之一的拉力对前进毫无贡献。
\end{example}

\begin{definition}[动力做功、阻力做功、零功]\label{def:g11:work-of-force:sign}
$\cos\theta$ 的符号把每个力分为三种情形：
\begin{itemize}
  \item $\theta < \qty{90}{\degree}$：$W > 0$，此时的功是
        \emph{动力功}\index{动力功}——力把能量输入
        运动之中；
  \item $\theta = \qty{90}{\degree}$：$W = 0$，力完全不做功，
        无论它有多大；
  \item $\theta > \qty{90}{\degree}$：$W < 0$，此时的功是
        \emph{阻力功}\index{阻力功}——力消耗
        能量。典型情形是 $\qty{180}{\degree}$ 的制动：
        \qty{800}{N} 的力作用\qty{45}{m}，使汽车失去 $800 \times 45 \times (-1)
        = \qty{-36}{kJ}$ 的能量，转化为制动盘中的热量。
\end{itemize}
\end{definition}

\begin{omfigure}
\begin{tikzpicture}[font=\small]
  \draw (0,0) -- (2.6,0); \draw[fill=omExo!15] (0.9,0) rectangle (1.7,0.55);
  \draw[-Stealth, omDef, very thick] (1.3,0.28) -- ++(40:1.1)
    node[above] {$\vect F$};
  \draw[-Stealth, omThm, thick] (0.5,-0.5) -- (2.1,-0.5);
  \node at (1.3,2.0) {$W > 0$};
\end{tikzpicture}\qquad
\begin{tikzpicture}[font=\small]
  \draw (0,0) -- (2.6,0); \draw[fill=omExo!15] (0.9,0) rectangle (1.7,0.55);
  \draw[-Stealth, omDef, very thick] (1.3,0.55) -- ++(90:1.1)
    node[above] {$\vect F$};
  \draw[-Stealth, omThm, thick] (0.5,-0.5) -- (2.1,-0.5);
  \node at (1.3,2.0) {$W = 0$};
\end{tikzpicture}\qquad
\begin{tikzpicture}[font=\small]
  \draw (0,0) -- (2.6,0); \draw[fill=omExo!15] (0.9,0) rectangle (1.7,0.55);
  \draw[-Stealth, omDef, very thick] (1.3,0.28) -- ++(140:1.1)
    node[above] {$\vect F$};
  \draw[-Stealth, omThm, thick] (0.5,-0.5) -- (2.1,-0.5);
  \node at (1.3,2.0) {$W < 0$};
\end{tikzpicture}

{\small 三种情形——动力、零功、阻力——由
力与位移（红色箭头）之间的夹角区分。}
\end{omfigure}

\section{重力的功}

\begin{proposition}[重力的功]\label{prop:g11:work-of-force:weight}
质量为 $m$ 的物体从 $A$（高度 $z_A$）运动到
$B$（高度 $z_B$）时，沿\emph{任意}由直线段组成的路径，
其重力 $\vect P$ 所做的功为
\[
W_{AB}(\vect P) = mg\,(z_A - z_B) = \pm\, mgh ,
\]
其中 $h = \abs{z_A - z_B}$ 为高度下降量：下降时为 $+mgh$，
上升时为 $-mgh$——功只取决于高度下降量，与路径无关。
\end{proposition}

\begin{proof}
在一段直线段 $AB$ 上，重力（大小为 $mg$，竖直向下）
$\theta$ 为重力与 $\vect{AB}$ 的夹角，而 $AB\cos\theta$ 恰好是
位移在竖直向下方向上的投影：
$z_A - z_B$。于是 $W = mg \times AB\cos\theta = mg\,(z_A - z_B)$。
沿一串线段累加时，各段的功相加，高度差依次相消：
$mg(z_A - z_C) + mg(z_C - z_B) = mg(z_A - z_B)$.
\end{proof}

\begin{proposition}[任意路径]\label{prop:g11:work-of-force:anypath}
从 $A$ 到 $B$ 的任意曲线路径上，同一公式同样成立。
\end{proposition}
\admitted

\begin{remark}\label{rem:g11:work-of-force:anypath}
这一结论合理——曲线可以任意接近由许多短线段组成的折线，
但要把“任意逼近”讲清楚，
需要微积分，见\cref{ch:g12:work-and-energy}。
\end{remark}

\begin{omfigure}
\begin{tikzpicture}[font=\small]
  \coordinate (A) at (0.7,2.6); \coordinate (B) at (4.4,0);
  \draw[dashed, omExo] (0,2.6) -- (5.2,2.6);
  \draw[dashed, omExo] (0,0) -- (5.2,0);
  \draw[Stealth-Stealth, omExo] (5.0,0.05) -- (5.0,2.55)
    node[midway, right] {$h$};
  \draw[omDef, very thick] (A) -- (B);
  \draw[omMeth, very thick] (A) -- (1.5,2.0) -- (3.1,2.0) --
    (2.3,0.9) -- (3.6,0.55) -- (B);
  \draw[-Stealth, omThm, thick] (2.55,1.85) -- ++(0,-0.8)
    node[below] {$\vect P$};
  \fill (A) circle (1.5pt) node[above] {$A$};
  \fill (B) circle (1.5pt) node[below] {$B$};
\end{tikzpicture}

{\small 直线（蓝色）或曲折（绿色）的路径，从 $A$
到 $B$ 下降的高度 $h$ 都相同：重力做功均为 $+mgh$。}
\end{omfigure}

\begin{example}[登山者]\label{ex:g11:work-of-force:hiker}
一位携带背包的 \qty{78}{kg} 登山者，攀登 \qty{450}{m}；
无论沿哪条路线，$W(\vect P) = -78 \times 9.81 \times 450 \approx \qty{-344}{kJ}$，
下山时，无论绕行多少，重力做功恰为 $+\qty{344}{kJ}$。
\end{example}

\section{摩擦力做的功}

\begin{proposition}[滑动摩擦力做的功]\label{prop:g11:work-of-force:friction}
大小恒为 $f$ 的摩擦力，在每一时刻都与运动方向相反，
沿路径做功 $W = -fL$，其中 $L$ 为路径总长：摩擦力总是阻碍运动，
且与\emph{所经路径的长度}成正比，
而不与位移有关。
\end{proposition}

\begin{proof}
每一小段直线上，摩擦力与运动方向的夹角为 $\qty{180}{\degree}$，
贡献 $-f \times (\text{该段长度})$；各段相加得
$-fL$。曲线路径同样切成小段处理，
其极限在本阶段不加证明地承认，与重力的情形相同（见\cref{rem:g11:work-of-force:anypath}）。
\end{proof}

\begin{remark}[摩擦力与重力：分水岭]\label{rem:g11:work-of-force:divide}
重力做的功与路径无关；摩擦力做的功则如同每米收取的“过路费”。
绕一圈回到原处时，重力做的功净值为零——上坡时它所“取走”的，
下坡时又全部归还；而摩擦力每一段都要“收费”$-fL$，
却从不归还：重力做的功可以被储存并回收（势能，
见\cref{ch:g11:mechanical-energy}），摩擦力做的功则以热量的形式损失。
\end{remark}

\begin{example}[往返一次]\label{ex:g11:work-of-force:roundtrip}
一个箱子沿地面滑动 \qty{5.0}{m}，再滑回 \qty{5.0}{m}，
摩擦力为 $f = \qty{60}{N}$，故摩擦力做功 $-60 \times 10 =
\qty{-600}{J}$，尽管箱子最终回到了出发点。
重力与法向力都与运动方向垂直，全程做功为零。
\end{example}

\begin{remark}[弹簧呢？]\label{rem:g11:work-of-force:spring}
弹簧伸长得越多，拉力越大：其力并非恒力，
因此\cref{def:g11:work-of-force:work}不再适用。
求弹簧做的功，需要把伸长量切成无穷小的步长，
再逐段求和——这是微积分的内容，见\cref{ch:g12:work-and-energy}。
\end{remark}

\section{功率}

\begin{definition}[功率]\label{def:g11:work-of-force:power}
力的\emph{功率}\index{功率!机械}，是指力做功的快慢，
$P = W/\Delta t$，单位为\emph{瓦特}\index{瓦特}
（$\unit{W} = \unit{J/s}$），与
\cref{ch:g11:circuits-and-power}中的瓦特是同一个单位。一个前国际单位制的单位仍保留在发动机铭牌上，
即\emph{马力}\index{马力}，
$\qty{1}{hp} = \qty{736}{W}$。
\end{definition}

\begin{proposition}[匀速时的功率]\label{prop:g11:work-of-force:pfv}
若物体以恒定速率 $v$ 运动，则
恒力 $\vect F$ 与运动方向成夹角 $\theta$，所提供的功率为
$P = F v \cos\theta$；特别地，
若力沿运动方向，则 $P = Fv$。
\end{proposition}

\begin{proof}
在时间 $\Delta t$ 内，物体移动的距离为 $d = v\,\Delta t$，故
$W = F \times v\,\Delta t \times \cos\theta$；两边除以 $\Delta t$。
\end{proof}

\begin{example}[起重机]\label{ex:g11:work-of-force:crane}
起重机吊起一个 \qty{600}{kg} 的托盘，速度恒为 \qty{0.90}{m/s}：
钢缆的拉力等于重力，$mg = \qty{5.89e3}{N}$，故
$P = Fv = 5890 \times 0.90 \approx \qty{5.3}{kW} \approx \qty{7.2}{hp}$
——相当于一支七匹马的车队，装在一只钢箱里。
\end{example}

\begin{method}[以较小的力爬坡：之字形路线]\label{met:g11:work-of-force:zigzag}
用有限的力将重物提升高度 $h$：
\begin{enumerate}
  \item 总功是固定的：由\cref{prop:g11:work-of-force:weight}，
        无论走哪条路到达顶端，都需做功 $mgh$。
  \item 拉长路径：在长为 $L$ 的斜坡上，所需的力
        （匀速，不计摩擦）为 $F = mgh/L$——路长加倍，
        力减半。
  \item 付出附加代价：摩擦力按路程累加计费
        （见\cref{prop:g11:work-of-force:friction}），额外增加 $fL$。
\end{enumerate}
山路之字形盘绕，正是为了第 2 步：高度相同，力更小，做功不变。
\end{method}

\begin{omfigure}
\begin{tikzpicture}
\begin{axis}[omaxis, width=8.8cm, height=5.0cm,
  xlabel={路径长度 $L$ (\unit{km})},
  ylabel={所需的力 (\unit{kN})},
  xmin=0.8, xmax=6.2, ymin=0, ymax=3.2,
  domain=1.2:6, samples=100]
  \addplot[omDef, thick] {3.53/x};
  \addplot[only marks, mark=*, omThm] coordinates {(1.5,2.35) (5,0.71)};
  \node[font=\small, anchor=west] at (axis cs:1.65,2.35) {直线轨道};
  \node[font=\small, anchor=south west] at (axis cs:4.0,0.8) {之字形山路};
\end{axis}
\end{tikzpicture}

{\small 将 \qty{1200}{kg} 的物体提升 $h = \qty{300}{m}$，其乘积
$F \times L = mgh = \qty{3.53}{MJ}$ 为定值，故山路以增加路长
换取较小的力（见\cref{exo:g11:work-of-force:12}）。}
\end{omfigure}

\section{练习}

\begin{exercise}[$\star$]\label{exo:g11:work-of-force:1}
恒力 $F = \qty{40}{N}$ 沿直线作用了 \qty{50}{m}
的位移，求当 $\theta = \qty{0}{\degree}$、
$\qty{60}{\degree}$、$\qty{90}{\degree}$ 和 $\qty{120}{\degree}$ 时该力所做的功。
\end{exercise}

\begin{exercise}[$\star$]\label{exo:g11:work-of-force:2}
一个 \qty{12}{kg} 的工具箱从货车上放下 \qty{2.5}{m}，之后
又被提回车上，求每种情况下重力所做的功？
\end{exercise}

\begin{exercise}[$\star$]\label{exo:g11:work-of-force:3}
判断下列各力做功为驱动功、阻力功还是零？(a) 下落苹果的重力；(b) 上升
小球的重力；(c) 滑动箱子所受的法向力；(d)
刹车自行车所受的摩擦力；(e) 摆动中摆球
所受细绳的张力。
\end{exercise}

\begin{exercise}[$\star$]\label{exo:g11:work-of-force:4}
一台起重机以匀速将 \qty{600}{kg} 的物体提升 \qty{25}{m}。求提升力
与重力所做的功？它们的和是多少？为什么？
\end{exercise}

\begin{exercise}[$\star$]\label{exo:g11:work-of-force:5}
一台自动扶梯的电动机每分钟做功 \qty{90}{kJ}。其功率
为多少瓦特？合多少马力？
\end{exercise}

\begin{exercise}[$\star\star$]\label{exo:g11:work-of-force:6}
一只雪橇被张力为 \qty{120}{N} 的绳子拖行 \qty{200}{m}，绳子
与雪面成 $\qty{30}{\degree}$ 斜向上；摩擦力为 \qty{45}{N}。求
四个力（绳的拉力、摩擦力、重力、法向力）各自所做的功及总功？
\end{exercise}

\begin{exercise}[$\star\star$]\label{exo:g11:work-of-force:7}
一名骑车人以恒定速率 \qty{25}{km/h} 骑行，克服总阻力
\qty{18}{N}。求其输出功率？一小时内做功多少？
\end{exercise}

\begin{exercise}[$\star\star$]\label{exo:g11:work-of-force:8}
一名 \qty{65}{kg} 的滑雪者下降 \qty{120}{m} 的高度，可沿 \qty{800}{m}
的陡坡或 \qty{1500}{m} 的缓道滑下；两条路线上滑雪板所受摩擦力均为
\qty{30}{N}。求每条路线上重力与摩擦力所做的功？
哪种功与路线有关？
\end{exercise}

\begin{exercise}[$\star\star$]\label{exo:g11:work-of-force:9}
一个 \qty{20}{kg} 的箱子沿倾角 $\qty{25}{\degree}$ 的斜面向上
推动 \qty{4.0}{m}，匀速运动，克服 \qty{30}{N} 的摩擦力。求
重力与摩擦力所做的功，然后由匀速条件求
平行于斜面的推力所做的功及其大小。
\end{exercise}

\begin{exercise}[$\star\star$]\label{exo:g11:work-of-force:10}
一辆汽车以 \qty{130}{km/h} 的速率巡航，发动机向车轮输出
\qty{50}{kW} 的功率。它所克服的总阻力为多大？
\end{exercise}

\begin{exercise}[$\star\star$]\label{exo:g11:work-of-force:11}
一匹纤马以 \qty{1.0}{m/s} 的速度拖着驳船，
张力为 \qty{800}{N} 的绳子与纤道成 $\qty{20}{\degree}$ 角。求
纤马的功率（分别以瓦特和马力表示），并加以评论。
\end{exercise}

\begin{exercise}[$\star\star\star$]\label{exo:g11:work-of-force:12}
一辆 \qty{1200}{kg} 的汽车爬升 $h = \qty{300}{m}$ 至山口，可走
坡度为 $20\%$ 的直线轨道，或坡度为 $6\%$ 的之字形山路（坡度 $=
\sin\alpha$）。忽略摩擦，匀速行驶，对每条路线求
路程、沿斜面所需的力及其所做的功。
用一句话作结论。
\end{exercise}

\begin{exercise}[$\star\star\star$]\label{exo:g11:work-of-force:13}
一台绞车将 \qty{250}{kg} 瓷砖沿脚手架向上提升 \qty{15}{m}，
以稳定速度 \qty{0.40}{m/s} 运动。求缆绳拉力所做的功、
所需时长、机械功率（两种方法：$W/\Delta t$ 与 $Fv$），以及
电动机效率为 $70\%$ 时消耗的电功率。
\end{exercise}

\begin{exercise}[$\star\star\star$]\label{exo:g11:work-of-force:14}
将一衣柜克服 $f = \qty{140}{N}$ 的摩擦力，从房间的一角
推至对角：可沿 \qty{6.0}{m} 的对角线直线移动，
或沿墙先后移动 \qty{3.6}{m} 与 \qty{4.8}{m}。求摩擦力在每条路线上所做的功，
并求沿对角线往返一次时摩擦力所做的功？重力在
该往返过程中做功多少？这一结果揭示了两种力之间
怎样的本质区别？
\end{exercise}

\begin{exercise}[$\star\star\star$]\label{exo:g11:work-of-force:15}
一名骑车者（连车共 \qty{80}{kg}）翻越一座山口：
路程 \qty{12}{km}，爬升高度 \qty{600}{m}，以恒定 \qty{15}{km/h}
速率骑行，并克服 \qty{14}{N} 的空气阻力与滚动阻力。求
克服重力所做的功、克服摩擦力所做的功及总功；
爬坡时长；以瓦特和马力分别表示的平均功率。
\end{exercise}

\section{问题：搬运工的一天}

\begin{problem}[{周末问题——物理以焦耳计价：一天的推移、吊运与
发夹弯，
最后谦卑地发现水壶对诚实劳动的看法}]
\label{pb:g11:work-of-force:1}
一班搬运工清空一座仓库，把货物吊运至三楼的
住宅，然后把货车开往山顶村庄。每项任务都以焦耳计价，
厨房的水壶将审核这一天的账目。每箱质量 \qty{30}{kg}。

\medskip
\noindent\textbf{第一部分——穿过仓库地面。}
一只箱子受水平推力 $F = \qty{110}{N}$ 的作用，
沿地面移动 $d = \qty{8.0}{m}$，滑动摩擦力为 $f = \qty{95}{N}$。
\begin{enumerate}
  \item 求推力所做的功。
  \item 求摩擦力、重力与
        法向力所做的功（并说明为零的理由）。
  \item 求箱子所受的总功，并说明其正负号的含义（箱子原来
        处于静止），什么样的推力能使速率保持不变
        （见\cref{ch:g11:forces-and-motion}）？
  \item 到达远端角落，可沿 \qty{10.0}{m} 的对角线，
        或沿两面墙（先 \qty{6.0}{m} 后 \qty{8.0}{m}）。摩擦力
        在每条路线上所做的功各为多少？
  \item 沿对角线往返（去程与回程），求
        摩擦力与重力所做的功。哪个力会“退还”功？
\end{enumerate}

\medskip
\noindent\textbf{第二部分——上到三楼（$h = \qty{9.0}{m}$）。}
\begin{enumerate}[resume]
  \item 求箱子被提升至住宅时重力所做的功。
        搬运工至少需要做多少功？
  \item 一名 \qty{70}{kg} 的搬运工带着箱子上楼，
        重力对搬运工\emph{和}箱子都“收费”：求克服重力的总功，
        以及其中真正用在箱子上的比例是多少？
  \item 手推车用的斜坡长 \qty{30}{m}，升高同样为 \qty{9.0}{m}，
        滚动摩擦力为 \qty{40}{N}。匀速运动时，
        求沿斜坡所需的推力及其所做的功。
  \item 比较斜坡与竖直吊运：所需的力、
        所做的功各为多少？斜坡换来了什么，又付出了什么代价？
  \item 一套绳索与滑轮装置以 \qty{0.50}{m/s} 的速度提升箱子。
        求所需时长与机械功率？
  \item 用 $P = Fv$ 重新求出该功率。吊装电动机的效率为 $60\%$，
        求其消耗的电功率？
\end{enumerate}

\medskip
\noindent\textbf{第三部分——盘山公路。}
满载的货车（\qty{3500}{kg}）须抵达比山谷高 \qty{240}{m} 的村庄，
有两条路线：直行车道长 \qty{1.0}{km}，坡度 $24\%$；或
铺装道路长 \qty{8.0}{km}，坡度 $3.0\%$（坡度 $= \sin\alpha$）。
\begin{enumerate}[resume]
  \item 验证两条路线的爬升高度均为 \qty{240}{m}。
  \item 忽略摩擦，匀速行驶时，求
        每条路线上所需的驱动力。
  \item 驱动力在每条路线上所做的功各为多少？与 $mgh$ 比较。
  \item 两条路线的滚动阻力均为 \qty{400}{N}，求
        每条路线上额外所做的功。盘山公路的舒适性要付出什么代价？
  \item 以 \qty{36}{km/h} 的速度行驶，求
        每条路线所需的发动机功率（分别以 \unit{kW} 和 \unit{hp} 表示）。道路为何要盘旋？
\end{enumerate}

\medskip
\noindent\textbf{第四部分——当日账目。}
\begin{enumerate}[resume]
  \item 搬运工将 $60$ 只箱子推过地面（即第 1 题中的
        推力），并把全部 $60$ 只箱子吊上 \qty{9.0}{m} 高处
        （即第 6 题）。输送给箱子的机械功总计为多少？
  \item 再计入上楼搬运工本人：$20$ 次上楼，人体质量为
        \qty{70}{kg}（下楼时的“返还”功归于发热的膝盖，不计入账目）。当日
        总功为多少？
  \item 搬运工工作了 $8.0$ 小时：平均机械功率为多少？
        一只 \qty{2.0}{kW} 的水壶需多长时间才能消耗同样的能量？
  \item 结语：用两句话说明焦耳计量的是什么——
        以及它不计量什么（墙壁被推了一下午，箱子却静止不动）
——并说明搬运工为何拥有绞车、斜坡与盘山
        公路，而不是更大的肌肉。
\end{enumerate}
\end{problem}
